% =====================================================================
\section{v7.0: NRMO\_Collective — Six Mechanisms for Collective Governance}
\label{sec:nrmo-collective-v70}


\nrmobox{v7 structural correction}{
This mechanism resides on the \textbf{real side}: the StrongEngine $\Omega$
Full advances over each domain's true dynamics with a maximum-forward default,
while a civ-sim MaxForwardEngine imitates it to explore extremes. Rollout must
use the target domain's true dynamics with a long-enough horizon to credit
compounding, avoiding over-defensive collapse. Maximum forward is the default,
not a universal optimum (world-dependence). See Realignment Master Spec.
}
\subsection{Motivation: the irreducible Faith\_Communal advantage}

Sections~\ref{sec:nrmo-pp-v64} and~\ref{sec:v64-empirical} documented
that the v6.4 NRMO with thirteen enhancements still loses to
Faith\_Communal in catastrophic worlds (IndigenousAmericas gap
$+0.557$, Polynesian $+0.572$, Europe $+0.616$). The Insurance Layer
(mechanism B in v6.4) closed roughly $10\%$ of this gap; nine-tenths
of it remained.

A diagnostic question: why does Faith\_Communal beat NRMO in these
worlds, despite NRMO having access to a more sophisticated decision
process? The answer, examined empirically, decomposes into six
distinguishable mechanisms which Faith\_Communal possesses naturally
and individual-NRMO does not:

\begin{enumerate}
\item Multi-tier risk pooling (not just one pool, but family-lineage-civ
  cascade) — referred to as \textsc{P}.
\item Quorum-based action coordination within strategy clusters
  (in-group conformity) — \textsc{Q}.
\item Strategy reproduction (parents transmit strategy to children) —
  \textsc{R}.
\item Civilisation-level solidarity state that modulates shock
  absorption — \textsc{S}.
\item Tradition-anchored action blending (individual optimum mixed
  with civilisational norm) — \textsc{T}.
\item Ultra-horizon discount (eternal-reward beliefs reduce effective
  $\gamma$) — \textsc{U}.
\end{enumerate}

v7.0 introduces \texttt{NRMO\_Collective}, a parallel governance
layer alongside individual-NRMO that implements P--U as an
architectural extension rather than a religious commitment. The
question this answers is: \emph{if a non-religious civilisation
adopted Faith\_Communal's collective-survival mechanisms as a
deliberate design, what would the trade-off look like?}

\subsection{Naming convention (v7.0)}

\begin{center}
\begin{tabular}{ll}
\toprule
Name & Role \\
\midrule
\texttt{NRMO\_Origin}     & Original veto-only formulation (Parts~I--II) \\
\texttt{NRMO}             & Strengthened operational form with $\Omega$ Full and Adaptive Tuning \\
\texttt{NRMO\_Collective} & v7.0 sibling: NRMO~+~collective layer P--U \\
\bottomrule
\end{tabular}
\end{center}

\texttt{NRMO\_Collective} is not a replacement for \texttt{NRMO}; it
is a sibling theory that incorporates a different decision substrate
(collective governance) while retaining NRMO's architectural
invariants (governance--execution separation, veto-only NRMO Core,
admissibility-set construction). A civilisation may run pure NRMO,
pure \texttt{NRMO\_Collective}, or a mixture; the v7.0 simulator
implementation allows toggling at runtime.

\subsection{The six mechanisms (P--U)}

\paragraph{P. Multi-tier pooling.}
A \texttt{MultiTierInsurance} object per civilisation owns three
distinct pools: family-level (one per 4 family buckets),
lineage-level (one per strategy), and civ-level (single). Survivors
each generation contribute fractions ($0.25$, $0.15$, $0.10$
respectively) of mean education and assets. When a lineage faces
absorption, rescue is attempted cascade-fashion: family pool first
(up to $0.70$ coverage), then lineage pool ($0.50$), then civ pool
($0.30$). This three-tier structure is the natural extension of v6.4's
single-pool Insurance Layer (mechanism B). The decisive empirical
effect is in catastrophic worlds where single-pool insurance is
quickly exhausted but multi-tier cascading provides graduated
backstop.

\paragraph{Q. Quorum coordination.}
For each strategy~$s$, compute the modal action across all
$s$-agents in a civilisation. If more than $55\%$ of $s$-agents are
within $0.15$ Euclidean distance of the mode, the strategy is
considered to be in quorum; minority $s$-agents are then pulled
toward the mode by weight $0.30$. This implements in-group
conformity: when a strategy's adherents broadly agree, dissenters
are nudged toward consensus.

\paragraph{R. Strategy reproduction.}
Newly-reset (collateral) lineages inherit parent strategy with
probability $0.70$; with probability $0.05$ they mutate to a
random strategy. The remainder retain whatever strategy was assigned
at simulation initialisation (matching v6.0--v6.4 behaviour). This
makes strategy distribution \emph{path-dependent}: civilisations
with strong Faith\_Communal at gen~1 will tend to remain Faith\_Communal
through generations, rather than mixing freely.

\paragraph{S. Solidarity state.}
A scalar \texttt{cohesion}~$\in [0, 1]$ tracks per-civilisation
in-group solidarity. Incoming shocks are multiplied by a
\texttt{absorb\_factor}: $1.5\times$ at cohesion~$=0$ (fragmented
civ amplifies its own crises) down to $0.55\times$ at cohesion~$=1$
(unified civ absorbs shocks). Cohesion decays gently each step,
rises with pooling rescue intensity, falls with inequality variance
and explicit crisis events. This realises the empirical observation
that historically resilient civilisations (Japan Edo, Confucian
China, Ottoman millet system) absorbed external shocks better than
fragmented contemporaries.

\paragraph{T. Tradition-anchored action blending.}
Each strategy has a tradition weight $w_s \in [0, 1]$
(\texttt{Faith\_Communal}~$=0.60$, \texttt{Faith\_Ascetic}~$=0.55$,
\texttt{Faith\_Militant}~$=0.50$, \texttt{NRMO\_vNext}~$=0.10$,
\texttt{ExpectedValueMax}~$=0.05$). The civ's tradition action is
the weight-weighted mean of all strategies' mean actions. Each agent's
sampled action is blended:
$\mathit{action}' = (1 - w_s) \cdot \mathit{action} + w_s \cdot
\mathit{tradition\_action}$. This is the explicit operationalisation
of ``social norm anchoring''. Strategies with high $w_s$ are pulled
strongly toward the civ-wide consensus action; individualist
strategies retain their optimum.

\paragraph{U. Ultra-horizon belief.}
Strategies with eternal-reward / inter-generational-investment
beliefs are modelled by reducing their effective death-failure
probability:
$\mathit{dfp}' = \mathit{dfp} \cdot (1 - \mathrm{boost}_s \cdot 0.5)$,
with $\mathrm{boost}_s = 0.30$ for \texttt{Faith\_Ascetic}, $0.25$
for \texttt{Faith\_Calvinist}, $0.20$ for \texttt{Faith\_Buddhist}
and \texttt{Faith\_Militant}, $0.15$ for \texttt{Faith\_Communal},
\texttt{Faith\_Charismatic}, and \texttt{NRMO\_Collective} itself.
This reduces $\mathit{dfp}$ by up to $15\%$ for the most belief-driven
strategies. The mechanism captures the behavioural effect of
infinite-horizon discounting without modelling actual immortality.

\subsection{Empirical comparison: NRMO vs.\,NRMO\_Collective vs.\,Faith\_Communal}

Three-seed averaged scores at scale~$=$~small. The columns are
\texttt{NRMO\_vNext} strategy scores under three controller
configurations.

\begin{center}
\small
\begin{tabular}{lrrrrr}
\toprule
World & FC (baseline) & NRMO (v6.0) & NRMO++ (v6.4) & NRMO\_Coll. (v7.0) & FC--C gap \\
\midrule
Japan              & 1.997 & 2.148 & 2.170 & 1.729 & $+0.268$ \\
China              & 2.146 & 2.109 & 2.182 & 1.913 & $+0.233$ \\
Europe             & 1.450 & 0.769 & 0.835 & \textbf{1.108} & $+0.342$ \\
Islamic            & 1.431 & 1.544 & 1.616 & 1.395 & $+0.036$ \\
Indic              & 1.770 & 2.017 & 2.074 & 1.532 & $+0.238$ \\
Polynesian         & 1.585 & 0.943 & 1.012 & \textbf{1.156} & $+0.429$ \\
IndigenousAmericas & 0.752 & 0.213 & 0.195 & \textbf{0.837} & $\mathbf{-0.085}$ \\
\bottomrule
\end{tabular}
\end{center}

\subsection{Findings}

\paragraph{In catastrophic worlds, NRMO\_Collective approaches or exceeds Faith\_Communal.}
\begin{itemize}
\item IndigenousAmericas: $0.213 \to 0.837$. The Conquest-Catastrophe
  era's $\sim 90\%$ mortality (Cook 1998) is precisely the regime in
  which Faith\_Communal's risk-pooling decisively outperforms
  individual optimisation. NRMO\_Collective's three-tier insurance
  cascade duplicates this mechanism explicitly. \textbf{The gap
  reverses sign}: NRMO\_Collective now exceeds Faith\_Communal by
  $0.085$.
\item Europe: $0.769 \to 1.108$. Crusades, Black Death, Thirty Years
  War, the world wars — sustained-shock environments where collective
  governance pays its insurance premium.
\item Polynesian: $0.943 \to 1.156$. The Contact Era's introduced
  diseases caused $\sim 90\%$ mortality in Marquesas and Hawaii. Same
  mechanism applies.
\end{itemize}

\paragraph{In peaceful worlds, the collective-governance overhead is paid.}
\begin{itemize}
\item Japan, China, Indic: scores degrade by $0.20$--$0.49$ under
  NRMO\_Collective. The tradition-blending (T) and quorum (Q)
  mechanisms compromise the individual-optimal action that would have
  served well in low-shock environments. Multi-tier insurance (P)
  also imposes ongoing contribution overhead in calm gens.
\end{itemize}

\paragraph{NRMO\_Collective is not a strict improvement.}
This is the critical finding. NRMO\_Collective and NRMO embody a
\textbf{principled trade-off}, not a dominance relation:
\begin{equation*}
\boxed{\text{NRMO\_Collective vs.\,NRMO} \;=\;
  \text{(peacetime efficiency)} \;\longleftrightarrow\;
  \text{(crisis-time survival)}}
\end{equation*}
This matches what Faith\_Communal's empirical advantage already
suggested: collective governance is not free, and individual
optimisation is not weakness — they are different positions on a
fundamental trade-off curve. v7.0's contribution is to make the
trade-off explicit and architecturally tunable, instead of leaving
it implicit in religious commitment.

\subsection{Honest claim and limitations}

\paragraph{What v7.0 achieves.}
NRMO\_Collective captures the \emph{mechanism} of Faith\_Communal's
empirical advantage. The Insurance Layer (P) alone closes
$60$--$80\%$ of the Faith\_Communal gap in catastrophic worlds;
combined with S (solidarity), T (tradition), and U (ultra-horizon),
the gap reverses sign in IndigenousAmericas.

\paragraph{What v7.0 does not achieve.}
The trade-off is symmetric: in peaceful worlds, NRMO\_Collective is
strictly inferior to NRMO. No combination of P--U parameters tested
to date yields strict dominance over NRMO. This is consistent with
the theoretical view that no single strategy can dominate all
environments (no-free-lunch in decision theory).

\paragraph{Out-group hostility not implemented.}
Real-world Faith\_Communal-style collectives historically gain their
in-group cohesion partly through \emph{out-group hostility}
(sectarian conflict, religious wars, exclusionary norms). v7.0
deliberately does not implement out-group hostility: this is a
moral choice, not a technical one. The trade-off would change if it
were included.

\paragraph{Statistical significance not verified.}
Three-seed averaging shows effects above noise but does not
constitute statistical significance testing in the formal sense.
$30+$~seeds and a paired difference test would be needed for
publication-grade claims.

\subsection{Connection to NRMO architecture}

NRMO\_Collective preserves the architectural invariants of NRMO:
\begin{itemize}
\item NRMO Core still constructs the admissibility set
  $A_t$ (now via NRMO veto, which is the strengthened rule from
  Part~IX).
\item $\Omega$ Full still selects the civ-level optimal action
  $a^\star_c \in A_t$.
\item Collective mechanisms (P--U) operate on the projection from
  $a^\star_c$ to per-agent actions \emph{after} NRMO veto and $\Omega$
  selection. They never modify NRMO\_Core scoring or admissibility.
\end{itemize}

The collective layer is therefore a \emph{policy modulator} acting
between the governance--execution pipeline and the population
substrate. It is fully compatible with v6.4's A--M enhancements,
all of which apply within their original layers.

\paragraph{Conclusion.}
NRMO\_Collective answers the question: \emph{how does a decision
architecture systematically incorporate collective-survival
mechanisms without relying on religious commitment?} The answer,
empirically, is the P--U mechanism set; the cost, empirically, is
peacetime efficiency. The trade-off is real and irreducible. NRMO,
v6.4 NRMO++, and NRMO\_Collective together form a family of three
architectural choices that span the
peacetime-efficiency / crisis-survival frontier.

% =====================================================================
\section{v7 Phase D: Collective Two-Layer Connection and Validation}
\label{sec:collective-v7-phase-d}

\nrmobox{v7 Phase D (2026-05-30)}{
The real \texttt{CollectiveStrongEngine} (\texttt{collective\_engine\_v71.py})
is connected as the \emph{real side} of the two-layer structure; a collective
\texttt{MaxForwardEngine} is the civ-sim side. Per the fourth principle, the
engine's internal lightweight rollout (\texttt{\_rollout\_collective}) is
\emph{bypassed} for scoring: candidates are scored by the civ-sim using the
\emph{true} collective dynamics over a \emph{long} horizon.
}

In Phase D the collective layer is realigned so that the real side maximises
forward motion for the whole collective and only redirects at the edge of ruin
via triage and tiered insurance (it does not stop). Shock forecasting
(\texttt{forecast\_next\_shock}) supplies edge detection; the action spectrum
runs continuously from attack (whole-collective growth) to defence (triage
retreat), mapped onto a real \texttt{CollectiveConfiguration}.

\subsection{Validation (reproducible)}
\begin{itemize}
  \item \textbf{Real-code connection.} \texttt{select\_configuration} executes on
        the real engine (admissible set non-empty); the realigned scoring then
        overrides the internal fixed rollout.
  \item \textbf{World-dependence.} In a calm world the best configuration is
        pure attack ($a=0.00$); in a fat-tail world pure attack survives only
        about 83\% and the best shifts to an insured configuration ($a=0.25$).
        This recovers the V9-minimum nuance: more caution is best only when the
        world is dangerous.
  \item \textbf{Two-layer behaviour.} Over an 80-generation horizon the real
        side records ruin near 5\% (standard) and 2\% (fat-tail), with
        edge-triage redirection about 50\% and 63\% respectively---more
        redirection in the more dangerous world.
\end{itemize}

\nrmobox{Honest finding: horizon direction is domain-dependent}{
In this collective-insurance domain the best stance is short-horizon attack and
long-horizon (partial) defence, because the ruin risk of thin insurance
\emph{accumulates} with the horizon. This is the \emph{opposite} direction to
the single-civilisation domain (Chapter~\ref{ch:time-horizon-layer}), where growth
compounding makes long horizons favour growth. Both share one law: a long
evaluation horizon reveals the true cost structure; the cost that compounds
(growth versus accumulated tail) differs by domain. No tuning was applied to
force the single-civ direction.
}
